\phantomsection
\part{線型空間}

\begin{abstract}
  このパートでは，線型空間について学ぶ．
  高校までに学んだ「ベクトル」は，「平面や空間内で向きと大きさを持った量」として説明されることが多かった．
  一方，線型代数学では，ベクトルを「ある条件を満たす集合の元」として抽象的に扱う．
  この考え方により，数ベクトルだけでなく，行列や関数なども同じ枠組みで議論できるようになる．
\end{abstract}

\phantomsection
\section{線型空間}

\phantomsection
\subsection{線型空間の定義}


線型代数学では，ベクトルとは必ずしも矢印である必要はない．
実際に，ある条件を満たす集合の元を「ベクトル」とすることができる．

そこでまず，ベクトルを抽象化した概念である
\textbf{線型空間}\index{せんけいくうかん@線型空間}
を定義する．

\begin{definition}{線型空間}{線型空間}
  $K$を体とする．
  集合$V$が次の条件を満たすとき，
  $V$を体$K$上の\textbf{線型空間}という．

  また，ここでは写像$ V \times V \to V$と$ K \times V \to V$を区別することなく\footnote{つまり，$\alpha + \beta$の$+$と$\alpha v + \beta w$の$+$を演算記号で区別しない．}演算記号$+$を用いると約束する．
\begin{enumerate}[label=(\roman*),leftmargin=3.2em]
  \item \textbf{加法に関して可換群} 
  
  集合$V$は，加法に関して可換群をなす．
  \item \textbf{スカラー倍に関する性質} 
 
  任意の$v,w\in V$および任意の$\alpha,\beta\in K$に対して，
  \[
  1_K v=v
  \]
  \[
  (\alpha+\beta) v=\alpha v+\beta v
  \]
  \[
  \alpha(v+w)=\alpha v+\alpha w,
  \]
  \[
  (\alpha\beta)v=\alpha(\beta v)
  \]
  が成り立つ．
  \end{enumerate}

  $V$の元を\textbf{ベクトル}という．
\end{definition}

\begin{mycolumn}
線型空間の公理のうち，
「加法に関して可換群である」という条件には，

\begin{itemize}
\item 結合法則
\item 交換法則
\item 零ベクトルの存在
\item 加法逆元の存在
\end{itemize}

が含まれている．具体的には，任意の$v,w,u\in V$に対して，
\[
(v+w)+u=v+(w+u),
\]
\[
v+w=w+v,
\]
\[
v+0_V=v,
\]
\[
v+(-v)=0_V
\]
が成り立つ．
\end{mycolumn}

\begin{example}{線型空間の例}{線型空間の例}
次の集合はいずれも線型空間となる．

\begin{itemize}
\item $\mathbb{R}^n$
\item $\mathbb{C}^n$
\item $m\times n$行列全体
\item 実数係数多項式全体
\end{itemize}

これらは一見異なる対象であるが，
線型空間という共通の枠組みで扱うことができる．
\end{example}

\subsection{線型部分空間}

線型空間の部分集合のうち，
それ自身も線型空間となるものを線型部分空間という．

\begin{definition}{線型部分空間}{線型部分空間}
  体$K$上の線型空間$V$の部分集合$W$が，
  次の条件を満たすとき，
  $W$を$V$の\textbf{線型部分空間}という．
\begin{enumerate}[label=(\roman*),leftmargin=3.2em]
    \item \textbf{空でない}
    \[
      W\neq\varnothing
    \]
    である．

    \item \textbf{加法に関する性質}

    任意の$v,w\in W$に対して，
    \[
      v+w\in W
    \]
    である．

    \item \textbf{スカラー倍に関する性質}

    任意の$v\in W$，
    任意の$\alpha\in K$に対して，
    \[
      \alpha v\in W
    \]
    である．
  \end{enumerate}
\end{definition}

\begin{mycolumn}
定義の最初の条件は
\[
0_V\in W
\]
としてもよい．

実際，
加法とスカラー倍について閉じている空でない集合は，
必ず零ベクトルを含むことが容易に示せる．
\end{mycolumn}

線型部分空間どうしから新たな線型部分空間を作る方法として，
共通部分と和が重要である．

\begin{definition}{共通部分・和}{共通部分・和}
  $V$を線型空間，
  $W_1,W_2$を$V$の線型部分空間とする．

  このとき，
  \[
  W_1\cap W_2
  \]
  を$W_1$と$W_2$の\textbf{共通部分}という．

  また，
  \[
  W_1+W_2
  =
  \{
  w_1+w_2
  \mid
  w_1\in W_1,\
  w_2\in W_2
  \}
  \]
  を$W_1$と$W_2$の\textbf{和}という．
\end{definition}

\begin{prop}{共通部分と和の線型部分空間性}{共通部分と和の線型部分空間性}
  \deref{共通部分・和}で定義した，
  線型部分空間の共通部分および和は，
  ともに線型部分空間である．
\end{prop}

\begin{leftbar}
\begin{proof}
  共通部分については，
  定義から容易に従う．

  また，
  和についても，
  加法およびスカラー倍について閉じていることを確かめればよい．
\end{proof}
\end{leftbar}

二つの線型空間を一つの線型空間として扱う方法として，
直和を定義する．

\begin{definition}{直和}{直和}
  $V,W$を体$K$上の線型空間とする．

  直積集合
  \[
  V\times W
  =
  \{
  (v,w)
  \mid
  v\in V,\
  w\in W
  \}
  \]
  を考え，
  加法およびスカラー倍を
  \[
  (v_1,w_1)+(v_2,w_2)
  =
  (v_1+v_2,w_1+w_2),
  \]
  \[
  \alpha(v,w)
  =
  (\alpha v,\alpha w)
  \]
  によって定める．

  このとき，
  $V\times W$は線型空間となる．

  これを
  $V$と$W$の\textbf{直和}といい，
  \[
  V\oplus W
  \]
  と表す．
\end{definition}

\begin{mycolumn}
直和
$V\oplus W$
は，
二つの線型空間を互いに独立な成分として並べた線型空間と考えることができる．
\end{mycolumn}
